\begin{proposition}[Subgroup criterion for finite groups] \label{proposition:subgroup_criterion_for_a_subset_of_a_finite_group}
    Let $G$ be a \CrefAndHyperrefIfExist{definition:group}{finite group} and let $H \subseteq G$ be a \CrefAndHyperrefIfExist{definition:subset_of_a_set}{subset}. Then $H$ is a \CrefAndHyperrefIfExist{definition:subgroup_of_a_group}{subgroup of $G$} if and only if the following conditions hold:
    \begin{enumerate}
        \item $H \neq \emptyset$,
        \item $H$ is closed under the group operation, i.e. For all $h_1, h_2 \in H$, the product $h_1 h_2$ is an element of $H$.
    \end{enumerate}
\end{proposition}
\begin{proof}
    By \Cref{proposition:subgroup_criterion_for_a_subset_of_a_group}, $H$ is a subgroup of $G$ if and only if the two conditions along with the following condition holds: that $H$ is closed under the inverse operator of $G$, i.e. for all $h \in H$, the inverse $h^{-1}$ is an element of $H$. Thus, it suffices to show that assuming that $G$ is a finite group, and the two conditions hold of $H$, then $H$ is also closed under the inverse operator of $G$. 


    For any $h \in H$, since $G$ is \CrefAndHyperrefIfExist{definition:finite_algebra_over_a_ring}{finite}, there exist some distinct positive integers $a,b$ such that $h^a = h^b$\CrefIfExists{definition:power_of_a_group_element_by_an_integer} --- if there did not, $G$ would have to have infinitely many distinct elements. Without loss of generality, say that $a > b$. In particular, $h^{a-b} = 1$\CrefIfExists{lemma:inverse_of_power_equals_power_of_inverse_in_a_group}\CrefIfExists{lemma:additivity_of_powers_in_a_group}, so $h^{a-b} \cdot h^{-1} = 1 \cdot h^{-1}$ and hence $h^{a-b-1} = h^{-1}$. Since $a > b$, the exponent $a-b-1$ is at least $0$. If it is $0$, then $1 = h^{-1}$, so $h^{-1}$ is an element of $H$. Otherwise, $a-b-1$ is positive, in which case $h^{a-b-1}$ can inductively be shown to belong to $H$ by the assumption that $H$ is closed under the group operation. Thus, $h^{-1}$ is an element of $H$. 
    \TODO{formulate the principle of mathematical induction in terms of $P(1)$ and $P(n)$ implies $P(n+1)$ for all positive integer $n$ implies that $P(n)$ holds for all positive integers $n$}
    \TODO{the induction that says that closure of group operation implies that closure this sorto f induction}
    \TODO{group element is identity iff inverse is identity}
\end{proof}